\chapter{CONCLUSION \& FUTURE WORK}
\section{Conclusion}
%  The conclusion paragraph should  summarize the Keep your conclusion positive.Emphasize the importance of your project Make sure the conclusion closes the project effectively from a reader's perspective.Reword and summarize the principal ideas of the project.
SecureEdgeNet successfully demonstrates a privacy-preserving Federated Learning based credit card fraud detection system using ensemble learning. The system trains Logistic Regression and MLP models across multiple simulated clients using Federated Averaging, while keeping transaction data local to each client partition. It also trains client-local XGBoost models to strengthen structured tabular fraud detection. The final fraud probability is generated by a weighted ensemble that combines global Logistic Regression, global MLP, and averaged XGBoost predictions.

The project addresses key challenges in fraud detection: class imbalance, privacy restrictions, model interpretability, distributed data ownership, and realistic evaluation. The system achieved strong test performance with an accuracy of 0.9634, precision of 0.7589, recall of 0.8673, F1-score of 0.8095, ROC-AUC of 0.9708, and PR-AUC of 0.8547. Since accuracy is misleading in fraud datasets, the recall, F1-score, PR-AUC, and confusion matrix provide more meaningful evidence of fraud detection quality. The modular design makes the project suitable for academic research, portfolio demonstration, and privacy-preserving machine learning education. It can run in Google Colab without complex infrastructure, while still demonstrating the essential concepts of Federated Learning, ensemble modelling, checkpointing, explainability, and inference.

\vspace{30mm}

\section{Future work}
%Future enhancements for a project can involve improving functionality, adding new features, or adapting to changing technologies. This could include refining the system to stay current with technical advancements, efficiently implementing new features, and integrating them with minimal modifications, ensuring the project remains flexible and adaptable. 
\begin{itemize}
    \item \textbf{Secure Aggregation:} Add cryptographic aggregation so the server can only observe summed updates, not individual client updates.
    \item \textbf{Differential Privacy:} Add update clipping and calibrated noise to provide formal privacy guarantees.
    \item \textbf{Federated XGBoost:} Replace local XGBoost ensembling with true federated gradient boosted tree protocols.
    \item \textbf{Non-IID Partitioning:} Simulate more realistic bank distributions using merchant, region, time, amount, or fraud-rate based partitions.
    \item \textbf{Graph Neural Networks:} Model relationships among cards, merchants, devices, IP addresses, and accounts to detect fraud rings.
    \item \textbf{Transformer Sequence Models:} Capture temporal transaction behaviour and long-range customer activity patterns.
    \item \textbf{Adaptive Ensemble Weights:} Learn or update ensemble weights based on validation performance and concept drift.
    \item \textbf{Production Security:} Add authentication, request validation, monitoring, and audit logs to the inference API.
\end{itemize}
